% Part: turing-machines
% Chapter: machines-computations
% Section: halting-states

\documentclass[../../../include/open-logic-section]{subfiles}

\begin{document}

\olfileid{tur}{mac}{hal}
\olsection{थांबण्याच्या अवस्था}

\begin{explain}
अमलात आणण्यासाठी कोणतीही सूचना उरली नसेल तेव्हाच आपली यंत्रे
थांबतात, अशी आपण व्याख्या केली आहे. पण ट्यूरिंग यंत्रांच्या
नेहमीच्या निरूपणांत थांबण्यासाठी खास \emph{थांबण्याची अवस्था}~$h$
असते; येथे $h \in Q$.

थांबण्याच्या अवस्थेमागील कल्पना सोपी आहे: यंत्राचे काम संपले की
(ते आदान स्वीकारण्यास तयार झाले किंवा त्याने प्रदान लिहून संपवले
की) ते~$h$ अवस्थेत जाते आणि तिथे थांबते. काही यंत्रांना अशा दोन
अवस्था असतात: एक आदान स्वीकारण्यासाठी आणि दुसरी ते नाकारण्यासाठी.
\end{explain}

\begin{ex}\emph{थांबण्याच्या अवस्था}.
ही कल्पना स्पष्ट करण्यासाठी सम यंत्रात बदल करून पाहू. आदानातील
रेघांची संख्या सम असेल तेव्हा यंत्र~$q_0$ अवस्थेत थांबण्याऐवजी,
त्याला थांबण्याच्या अवस्थेत पाठवणारी सूचना जोडू शकतो.
\[
\begin{tikzpicture}[->,>=stealth',shorten >=1pt,auto,node distance=2.8cm,
                    semithick]
  \tikzstyle{every state}=[fill=none,draw=black,text=black]

  \node[initial, state]         (A)                     {$q_0$};
  \node[state]         (B) [right of=A] {$q_1$};
  \node[state]         (C) [below of=A] {$h$};

  \path (A) edge [bend left] node {\TMtrans{\TMstroke}{\TMstroke}{R}} (B)
  	    edge node {\TMtrans{\TMblank}{\TMblank}{N}} (C)
        (B) edge [loop above] node {\TMtrans{\TMblank}{\TMblank}{R}} (B)
            edge [bend left] node {\TMtrans{\TMstroke}{\TMstroke}{R}} (A);
\end{tikzpicture}
\]

उदाहरण आणखी वाढवू. यंत्राला आदानातील रेघांची संख्या विषम असल्याचे
आढळल्यास मूळ सम यंत्र कधीच थांबत नाही. चक्रात फिरत राहणाऱ्या
सूचनेच्या जागी यंत्राला \emph{नकार} अवस्था~$r$ मध्ये नेणारी सूचना
देऊन यंत्रात बदल करू शकतो.
\[
\begin{tikzpicture}[->,>=stealth',shorten >=1pt,auto,node distance=2.8cm,
                    semithick]
  \tikzstyle{every state}=[fill=none,draw=black,text=black]

  \node[initial,state]         (A)                     {$q_0$};
  \node[state]         (B) [right of=A] {$q_1$};
  \node[state]         (C) [below of=A] {$h$};
  \node[state]         (D) [below of=B] {$r$};

  \path (A) edge [bend left] node {\TMtrans{\TMstroke}{\TMstroke}{R}} (B)
  	    edge node {\TMtrans{\TMblank}{\TMblank}{N}} (C)
        (B) edge node {\TMtrans{\TMblank}{\TMblank}{N}} (D)
            edge [bend left] node {\TMtrans{\TMstroke}{\TMstroke}{R}} (A);
\end{tikzpicture}
\]
\end{ex}

\begin{explain}
अशा प्रसंगी खास थांबण्याची अवस्था जोडल्याने यंत्र कोणती आदाने
स्वीकारते आणि कोणती नाकारते हे स्पष्ट होते. पण खास थांबण्याची
अवस्था जोडल्याने यंत्राची संगणनशक्ती वाढत नाही, हेही लक्षात
घ्यावे. त्याचप्रमाणे थांबण्याची कमी औपचारिक कल्पनाही उपयोगाची
आहे. या प्रकरणात आतापर्यंत वापरलेल्या थांबण्याच्या व्याख्येमुळे
\emph{थांबण्याची समस्या} याचा पुरावा अंतर्ज्ञानी रीतीने
सहज दाखवता येतो. म्हणून आपण आपली मूळ व्याख्या पुढेही वापरू.
\end{explain}

\end{document}
